\documentclass[11pt]{amsart}
\usepackage[T1]{fontenc}
\usepackage{amsmath,amssymb,amsthm}
\usepackage{hyperref}

\newtheorem{theorem}{Theorem}
\newtheorem{lemma}[theorem]{Lemma}
\theoremstyle{definition}
\newtheorem{definition}[theorem]{Definition}
\theoremstyle{remark}
\newtheorem{remark}[theorem]{Remark}

\newcommand{\K}{\mathbb K}
\newcommand{\Lcal}{\mathcal L}
\newcommand{\Zcal}{\mathcal Z}
\newcommand{\PiX}{\Pi(X)}

\title{Finite-Dimensional Spaces with Zero Numerical Index Have \(\mathbf L_{p,p}\)-nu}
\author{agent\_lane\_19}
\date{June 19, 2026}

\begin{document}
\maketitle

\section*{Source Problem}

The source paper is Sheldon Dantas, Sun Kwang Kim, Han Ju Lee, and Martin
Mazzitelli, \emph{On various types of density of numerical radius attaining
operators}, arXiv:2010.00280. After proving that finite-dimensional Banach
spaces with positive numerical index have \(\mathbf L_{p,p}\)-nu, and hence
that all finite-dimensional complex Banach spaces have it, the authors ask
whether the remaining real zero-index finite-dimensional case also has the
property:

\begin{quote}
Finite dimensional spaces with \(n(X)=0\) have the \(\mathbf L_{p,p}\)-nu?
\end{quote}

We answer this question affirmatively.

\section*{Definitions}

Let \(X\) be a real or complex Banach space. For \(T\in\Lcal(X)\), its
numerical radius is
\[
  v(T)=\sup\{|x^*(Tx)|:(x,x^*)\in\PiX\},
\]
where
\[
  \PiX=\{(x,x^*)\in S_X\times S_{X^*}:x^*(x)=1\}.
\]

The space \(X\) has \(\mathbf L_{p,p}\)-nu if for every \(\varepsilon>0\) and
every state \((x,x^*)\in\PiX\), there is
\(\eta(\varepsilon,(x,x^*))>0\) such that whenever \(T\in\Lcal(X)\) satisfies
\(v(T)=1\) and
\[
  |x^*(Tx)|>1-\eta(\varepsilon,(x,x^*)),
\]
there exists \(S\in\Lcal(X)\) such that \(v(S)=1\),
\[
  |x^*(Sx)|=1,
  \qquad
  \|S-T\|<\varepsilon .
\]

\section*{Idea of the Proof}

The positive-index proof in the source paper uses compactness of the
\(v\)-unit sphere in \(\Lcal(X)\). When \(n(X)=0\), \(v\) is only a seminorm,
so that sphere need not be bounded in the operator norm. The remedy is to
remove exactly the directions invisible to numerical radius.

Let \(\Zcal(X)=\{T\in\Lcal(X):v(T)=0\}\). On the quotient
\(\Lcal(X)/\Zcal(X)\), the numerical radius is a genuine norm. Since \(X\) is
finite-dimensional, this quotient is finite-dimensional, and the \(v\)-unit
sphere is compact. The state evaluation \(T\mapsto x^*(Tx)\) is constant on
cosets of \(\Zcal(X)\), so exact state attainment is a closed subset of that
compact sphere. Compactness gives a quotient-near exact attainer; lifting that
quotient perturbation gives an operator-norm-near exact attainer in
\(\Lcal(X)\).

\begin{lemma}\label{lem:quotient}
Let \(X\) be finite-dimensional and set
\[
  \Zcal(X)=\{T\in\Lcal(X):v(T)=0\}.
\]
For a fixed \((x,x^*)\in\PiX\), the map
\[
  \Phi_{x,x^*}(T+\Zcal(X))=x^*(Tx)
\]
is well-defined and continuous on \(\Lcal(X)/\Zcal(X)\). Moreover
\(\nu(T+\Zcal(X))=v(T)\) is a norm on the quotient.
\end{lemma}

\begin{proof}
If \(K\in\Zcal(X)\), then \(v(K)=0\), so
\[
  |y^*(Ky)|=0 \qquad ((y,y^*)\in\Pi(X)).
\]
In particular \(x^*(Kx)=0\), proving that \(\Phi_{x,x^*}\) is well-defined.
It is linear, hence continuous because the quotient is finite-dimensional.

The formula \(\nu(T+\Zcal(X))=v(T)\) is well-defined by the triangle
inequality for \(v\), and its kernel is exactly the zero coset by definition
of \(\Zcal(X)\). Thus \(\nu\) is a norm.
\end{proof}

\begin{theorem}\label{thm:main}
Every finite-dimensional real or complex Banach space has
\(\mathbf L_{p,p}\)-nu.
\end{theorem}

\begin{proof}
Fix \(\varepsilon>0\) and \((x,x^*)\in\PiX\). Let
\[
  Q=\Lcal(X)/\Zcal(X)
\]
and let \(q:\Lcal(X)\to Q\) be the quotient map. We use two norms on \(Q\):
\[
  \nu(q(T))=v(T),
  \qquad
  \|q(T)\|_q=\inf_{K\in\Zcal(X)}\|T-K\|.
\]
Since \(Q\) is finite-dimensional, these norms induce the same topology.

Consider the compact \(\nu\)-unit sphere
\[
  A=\{u\in Q:\nu(u)=1\}
\]
and the continuous function
\[
  f(u)=|\Phi_{x,x^*}(u)|.
\]
The closed set
\[
  M=\{u\in A:f(u)=1\}
\]
is nonempty: for \(R_0 z=x^*(z)x\), one has \(v(R_0)=1\) and
\(|x^*(R_0x)|=1\).

We claim that there is \(\eta>0\) such that every \(u\in A\) with
\(f(u)>1-\eta\) satisfies
\[
  \operatorname{dist}_q(u,M)<\varepsilon,
\]
where the distance is computed using the quotient operator norm
\(\|\cdot\|_q\). If not, there are \(u_n\in A\) with \(f(u_n)\to1\) and
\(\operatorname{dist}_q(u_n,M)\geq\varepsilon\). By compactness of \(A\), pass
to a subsequence \(u_n\to u_0\in A\). Continuity gives \(f(u_0)=1\), hence
\(u_0\in M\), contradicting the distance lower bound.

Now let \(T\in\Lcal(X)\) satisfy \(v(T)=1\) and
\[
  |x^*(Tx)|>1-\eta.
\]
Then \(q(T)\in A\) and \(f(q(T))>1-\eta\), so choose \(m=q(R)\in M\) and
\(K\in\Zcal(X)\) such that
\[
  \|T-(R+K)\|<\varepsilon.
\]
Put \(S=R+K\). Since \(q(S)=m\in M\), we have
\[
  v(S)=\nu(m)=1,
  \qquad
  |x^*(Sx)|=f(m)=1,
\]
and the displayed inequality gives \(\|S-T\|<\varepsilon\). This is exactly
\(\mathbf L_{p,p}\)-nu.
\end{proof}

\section*{Consequences and Scope}

The source paper already proved the theorem when \(n(X)>0\). Theorem
\ref{thm:main} removes the positivity assumption and answers the stated
finite-dimensional \(n(X)=0\) question affirmatively. It does not address the
other open questions in the source paper, such as whether
\(\mathbf L_{p,p}\)-nu implies strong subdifferentiability in the infinite
dimensional case, or whether \(\mathbf L_{o,o}\)-nu forces finite
dimensionality without additional assumptions.

\section*{Novelty Check}

I checked the run indexes for arXiv:2010.00280 and for the phrases
\(\mathbf L_{p,p}\)-nu, finite-dimensional \(n(X)=0\), zero numerical radius,
and quotient by skew-hermitian operators. No duplicate packet appeared. A web
search for the exact finite-dimensional zero-index question and close variants
found the source arXiv page and no separate later answer. The proof uses only
the quotient by zero numerical-radius operators already introduced in the
source paper.

\begin{thebibliography}{9}

\bibitem{DKLM}
S. Dantas, S. K. Kim, H. J. Lee, and M. Mazzitelli,
\emph{On various types of density of numerical radius attaining operators},
arXiv:2010.00280.

\end{thebibliography}

\end{document}
